\documentclass[12pt,a4paper]{article}

% --- Encodage & langue ---
\usepackage[utf8]{inputenc}
\usepackage[T1]{fontenc}
\usepackage[french]{babel}

% --- Mise en page ---
\usepackage[top=2.5cm, bottom=2.5cm, left=2.5cm, right=2.5cm]{geometry}
\usepackage{setspace}
\onehalfspacing
\usepackage{parskip}
\setlength{\parindent}{0pt}

% --- Typographie ---
\usepackage{lmodern}
\usepackage{microtype}

% --- Couleurs & graphiques ---
\usepackage[dvipsnames]{xcolor}
\usepackage{graphicx}
\usepackage{float}
\graphicspath{{../screenshots/}}
\usepackage{tikz}
\usetikzlibrary{shapes.geometric, arrows.meta, positioning, fit, backgrounds, calc}
\usepackage{tcolorbox}
\tcbuselibrary{skins, breakable}

% --- Tableaux ---
\usepackage{booktabs}
\usepackage{tabularx}
\usepackage{array}
\usepackage{multirow}
\usepackage{longtable}

% --- Code ---
\usepackage{listings}
\usepackage{fancyvrb}

\definecolor{codebg}{RGB}{248,248,248}
\definecolor{codeframe}{RGB}{200,200,200}
\definecolor{keywordcolor}{RGB}{0,0,180}
\definecolor{commentcolor}{RGB}{100,150,100}
\definecolor{stringcolor}{RGB}{180,0,0}

\lstset{
  language=Python,
  backgroundcolor=\color{codebg},
  frame=single,
  rulecolor=\color{codeframe},
  basicstyle=\ttfamily\footnotesize,
  keywordstyle=\color{keywordcolor}\bfseries,
  commentstyle=\color{commentcolor}\itshape,
  stringstyle=\color{stringcolor},
  showstringspaces=false,
  breaklines=true,
  numbers=left,
  numberstyle=\tiny\color{gray},
  numbersep=5pt,
  tabsize=4,
  captionpos=b,
}

% --- Liens & PDF ---
\usepackage[hidelinks]{hyperref}
\usepackage{nameref}

% --- En-têtes & pieds de page ---
\usepackage{fancyhdr}
\setlength{\headheight}{13.6pt}
\addtolength{\topmargin}{-1.6pt}
\pagestyle{fancy}
\fancyhf{}
\renewcommand{\headrulewidth}{0.4pt}
\fancyhead[L]{\small Projet de 1\iere{} année --- Groupe 53}
\fancyhead[R]{\small ENSAI 2025--2026}
\fancyfoot[C]{\thepage}

% --- Titres de sections ---
\usepackage{titlesec}
\titleformat{\section}{\large\bfseries\color{NavyBlue}}{}{0em}{\thesection\quad}
\titleformat{\subsection}{\normalsize\bfseries\color{NavyBlue!80}}{}{0em}{\thesubsection\quad}
\titleformat{\subsubsection}{\normalsize\itshape\bfseries}{}{0em}{\thesubsubsection\quad}

% --- Macros utilitaires ---
\newcommand{\sport}[1]{\textbf{#1}}
\newcommand{\code}[1]{\texttt{#1}}
\newcommand{\fichier}[1]{\texttt{#1}}
\newcommand{\classe}[1]{\texttt{\textbf{#1}}}

% Boîte de capture d'écran
\newcommand{\screenshot}[3]{%
  \begin{figure}[H]
    \centering
    \includegraphics[width=0.88\textwidth, height=0.45\textheight, keepaspectratio]{#3}
    \caption{#1}
    \label{fig:#3}
  \end{figure}
}

% Boîte de scénario
\newtcolorbox{scenario}[1][]{
  colback=NavyBlue!5!white,
  colframe=NavyBlue!40!white,
  fonttitle=\bfseries\small,
  title={#1},
  breakable,
  left=4pt, right=4pt, top=4pt, bottom=4pt
}

% Boîte de bilan personnel
\newtcolorbox{bilanbox}[1][]{
  colback=gray!5,
  colframe=gray!50,
  fonttitle=\bfseries,
  title={Bilan personnel},
  breakable
}

\begin{document}

% ============================================================
% PAGE DE GARDE
% ============================================================
\begin{titlepage}
  \centering
  \vspace*{1.5cm}

  {\LARGE\bfseries École nationale de la statistique\\
  et de l'analyse de l'information}\\[0.3cm]
  {\large ENSAI --- Bruz}\\[0.8cm]

  {\normalsize Projet traitement de données}\\[0.2cm]
  {\normalsize 1\iere{} Année --- 2025--2026}\\[1cm]

  \rule{\linewidth}{1.5pt}\\[0.4cm]
  {\Huge\bfseries\color{NavyBlue} Rapport du projet de\\traitement de données}\\[0.4cm]
  \rule{\linewidth}{1.5pt}\\[1.2cm]

  {\large\itshape Application en ligne de commande\\
  pour l'exploration de données sportives --- Python}\\[2cm]

  \textbf{Auteurs :}\\[0.4cm]
  \begin{tabular}{c}
    Mohammed-Adam \textsc{Lakhrissi} \\[0.15cm]
    Jade \textsc{Colin} \\[0.15cm]
    Mathias \textsc{Douillard} \\[0.15cm]
    Mohamed Fouad \textsc{Chaabane} \\
  \end{tabular}

  \vfill

  {\small \textit{Coordinateur :} Clément \textsc{Valot}\\[0.3cm]
  Projet de traitement de données --- 1\iere{} année\\
  Mai 2026}
\end{titlepage}

% ============================================================
% RÉSUMÉ
% ============================================================
\thispagestyle{empty}
\section*{Résumé}

Ce rapport présente le fonctionnement d'une application en ligne de commande (CLI)
développée en Python, permettant d'explorer et d'analyser des données statistiques issues
de neuf disciplines sportives et e-sportives : Basketball (NBA), Football (UEFA Champions
League et ligues européennes), Tennis (ATP/WTA), League of Legends (EMEA), Counter-Strike~2,
Échecs (FIDE), Volleyball (JO Paris 2024), StarCraft~II et Badminton (BWF).

L'application repose sur une architecture modulaire organisée en quatre couches : les
\textit{Parsers} chargent et normalisent les fichiers CSV, le \textit{Model} fournit des
classes génériques (\classe{Joueur}, \classe{Equipe}, \classe{Match}, \classe{Sport})
paramétrables pour chaque sport, les modules d'\textit{Analysis} implémentent les
fonctionnalités statistiques spécifiques, et le \textit{CLI} orchestre la navigation,
les filtres et l'export CSV.

Le projet totalise environ \textbf{11\,400 lignes de code Python} réparties sur
\textbf{69 fichiers}, exploite \textbf{36 fichiers CSV} de données couvrant plus de
\textbf{3\,000 joueurs} et \textbf{30\,000 matchs}, et dispose d'une suite de près de
\textbf{3\,800 lignes de tests unitaires} utilisant pytest. Après avoir présenté le
cahier des charges et les choix d'architecture, nous détaillerons la gestion des données,
la structure du code, l'utilisation de la programmation orientée objet et la répartition
des rôles, avant de présenter les résultats obtenus.

\bigskip
\textbf{Mots-clés :} Python, pandas, traitement de données, sports, CLI, programmation
orientée objet, tests unitaires.

\newpage

% ============================================================
% TABLE DES MATIÈRES
% ============================================================
\thispagestyle{empty}
\tableofcontents
\newpage

% ============================================================
\setcounter{page}{1}

% ------------------------------------------------------------
\section{Introduction}
% ------------------------------------------------------------

Suite à l'engouement suscité par les Jeux Olympiques et Paralympiques de Paris 2024,
le Ministère des Sports, de la Jeunesse et de la Vie Associative a exprimé le besoin de
se doter d'un outil informatique performant pour centraliser et gérer les résultats de
diverses compétitions sportives.

L'objectif principal de notre travail a été de développer une application générique,
permettant de modéliser, d'organiser et d'exploiter des entités universelles au monde du
sport : des joueurs, des équipes, des matchs et des compétitions. Cette application permet
à des utilisateurs d'accéder aux statistiques de ces compétitions, de les manipuler et d'en
déduire des analyses complexes.

La complexité de cette mission réside dans la diversité des disciplines traitées. Notre
application ne se limite pas à un seul sport, mais a été conçue pour intégrer aussi bien
des sports collectifs classiques comme le football ou le basketball, que des disciplines
individuelles comme le tennis ou les échecs, et même l'e-sport (League of Legends,
Counter-Strike~2, StarCraft~II). Chaque sport apportant ses propres spécificités
statistiques, l'enjeu de ce projet est de taille : concevoir une application capable
d'absorber des données hétérogènes tout en restant générique et évolutive, afin de
pouvoir intégrer facilement de nouveaux sports ou de nouvelles sources sans remettre en
cause son fonctionnement global.

Pour répondre à ces objectifs, nous avons adopté une démarche de \textbf{programmation
orientée objet}, complétée par l'utilisation de la bibliothèque \textbf{pandas} pour
optimiser le traitement de volumes de données importants. Nous avons construit nos
données comme des objets pour garantir une modélisation fidèle du domaine sportif,
tout en utilisant pandas pour assurer des performances élevées lors des calculs
statistiques et la gestion des valeurs manquantes. Nous justifierons ces choix
d'architecture et de programmation en détail dans la suite de ce document.

Les données proviennent de sources publiques variées (NBA API, ATP/WTA, Liquipedia,
FIDE, données des JO Paris 2024, etc.), toutes normalisées au format CSV. Le projet a
été développé collaborativement via Git et GitHub, en suivant une organisation modulaire
qui facilite l'ajout de nouveaux sports.

% ------------------------------------------------------------
\section{Cahier des charges}
\label{sec:cdc}
% ------------------------------------------------------------

L'application devait répondre à plusieurs exigences : couvrir au moins cinq sports
différents, proposer des fonctionnalités d'analyse statistique pertinentes pour chaque
sport, offrir une interface de navigation intuitive en ligne de commande, et être
accompagnée de tests unitaires (pytest, black, Flake8, docstrings NumPy). Nous avons
choisi de couvrir neuf disciplines pour enrichir la comparaison entre sports traditionnels
et e-sports.

\subsection{Diagramme de cas d'utilisation}

\begin{figure}[h!]
\centering
\begin{tikzpicture}[
  usecase/.style={draw, ellipse, fill=NavyBlue!8, minimum width=4.2cm,
                  minimum height=0.8cm, font=\small, align=center},
  sys/.style={draw, rectangle, dashed, rounded corners=4pt,
              minimum width=8cm, minimum height=11cm},
  arr/.style={->, thick, >=Stealth},
  inc/.style={->, dashed, thick, >=Stealth},
  >=Stealth
]
  % Boîte système
  \node[sys, anchor=center] (sys) at (2.5, 0) {};
  \node[above=5.2cm of sys.center, font=\bfseries\small] {Application Statistiques Sportives};

  % Cas d'utilisation
  \node[usecase] (filtrer)    at (2.5,  4.2) {Filtrer les sports};
  \node[usecase] (selec)      at (2.5,  2.8) {Sélectionner un sport};
  \node[usecase] (clast)      at (2.5,  1.5) {Voir le classement};
  \node[usecase] (stats)      at (2.5,  0.2) {Voir les stats d'une équipe};
  \node[usecase] (fiche)      at (2.5, -1.1) {Fiche individuelle joueur};
  \node[usecase, fill=orange!10] (roster) at (2.5, -2.4) {Voir le roster (collectif)};
  \node[usecase] (export)     at (2.5, -3.7) {Exporter en CSV};

  % Acteur
  \node[font=\small] (actlabel) at (-3.2, -4.5) {Utilisateur};
  \node[draw, circle, minimum size=0.6cm] (head) at (-3.2, 0.5) {};
  \draw (-3.2, 0.2) -- (-3.2, -0.8);
  \draw (-3.5, -0.2) -- (-2.9, -0.2);
  \draw (-3.2, -0.8) -- (-3.5, -1.4);
  \draw (-3.2, -0.8) -- (-2.9, -1.4);

  % Flèches acteur → cas
  \foreach \uc in {filtrer, selec, clast, stats, fiche, roster, export}
    \draw[arr] (-3.2, 0.2) -- (\uc.west);

  % Relations include
  \draw[inc] (filtrer.south) -- (selec.north)
    node[midway, right, font=\scriptsize\itshape] {«include»};
  \draw[inc] (selec.south) -- (clast.north)
    node[midway, right, font=\scriptsize\itshape] {«include»};

  % Relation extend
  \draw[inc] (fiche.east) to[out=0, in=0, looseness=1.5] (roster.east)
    node[midway, right=0.5cm, font=\scriptsize\itshape, align=left]
    {«extend»\\{[sport collectif]}};
\end{tikzpicture}
\caption{Diagramme de cas d'utilisation.}
\label{fig:usecase}
\end{figure}

Le diagramme de cas d'utilisation décrit les interactions entre l'utilisateur et
l'application. L'utilisateur dispose de deux actions de navigation. Il peut d'abord
\textbf{filtrer les sports} affichés en appliquant des critères comme le type de sport
(individuel ou collectif), la catégorie (sport traditionnel ou e-sport) ou le format de
compétition. Ce filtre inclut automatiquement l'étape suivante : la \textbf{sélection d'un
sport} dans la liste résultante.

Une fois un sport sélectionné, l'utilisateur accède aux fonctionnalités d'analyse : le
\textbf{classement}, les \textbf{statistiques d'équipe} ou la \textbf{fiche individuelle}
d'un joueur. La fonctionnalité \textit{voir le roster} est disponible uniquement pour les
sports collectifs (basketball, LoL, CS2, football, volleyball) --- c'est ce que représente
la relation \textbf{«extend»} conditionnée à [sport collectif]. Les relations
\textbf{«include»} indiquent des enchaînements obligatoires : filtrer les sports mène
systématiquement à en sélectionner un, et sélectionner un sport donne nécessairement
accès aux fonctionnalités d'analyse.

\subsection{Diagramme de classes}

Le diagramme ci-dessous présente l'architecture orientée objet du projet. Notre projet
est structuré en trois couches principales.

La première est la couche \textbf{Model}, qui contient les classes représentant les données
métier. On y trouve trois entités : \classe{Joueur}, \classe{Match} et \classe{Equipe}.
Chacune existe en double : une classe pour un objet individuel et une classe collection
(par exemple \classe{Joueurs}) qui regroupe tous les objets du même type et propose
des méthodes de filtrage et de statistiques.

Ces entités sont liées entre elles : une \classe{Equipe} peut récupérer la liste de ses
joueurs à partir d'une collection \classe{Joueurs}, et une collection \classe{Matchs}
permet de calculer le bilan victoires/défaites d'une équipe donnée. La classe \classe{Sport}
est indépendante des autres : elle décrit simplement les caractéristiques d'un sport via
trois énumérations, sans lien direct avec les joueurs ou les matchs.

La deuxième couche est \textbf{Parsers}. Elle contient la classe abstraite
\classe{BaseLoader} dont héritent les dix loaders, un par sport. Leur rôle est unique :
charger les fichiers CSV et retourner des données exploitables par le reste de
l'application.

La troisième couche est \textbf{Analysis}. Elle regroupe les modules de calcul statistique
propres à chaque sport. Ces modules s'appuient sur les loaders pour charger les données
et sur les classes du Model pour les manipuler. C'est la classe \classe{CLI} qui orchestre
l'ensemble en appelant ces modules selon les choix de l'utilisateur.

\subsection{Sports couverts et classification}

Chaque sport est caractérisé par trois axes de classification, modélisés par des enums
Python (\code{TypeSport}, \code{CategorieSport}, \code{TypeCompetition}) :

\begin{center}
\begin{tabular}{lccc}
  \toprule
  \textbf{Sport} & \textbf{Type} & \textbf{Catégorie} & \textbf{Compétition} \\
  \midrule
  Basketball (NBA)       & Collectif  & Sport   & Points       \\
  League of Legends      & Collectif  & E-sport & Points       \\
  Football Champions L.  & Collectif  & Sport   & Mixte        \\
  Football (ligues EU)   & Collectif  & Sport   & Points       \\
  Volleyball             & Collectif  & Sport   & Mixte        \\
  CS2                    & Collectif  & E-sport & Mixte        \\
  Tennis (ATP/WTA)       & Individuel & Sport   & Éliminatoire \\
  Badminton (BWF)        & Individuel & Sport   & Éliminatoire \\
  Échecs (FIDE)          & Individuel & E-sport & Éliminatoire \\
  StarCraft~II           & Individuel & E-sport & Éliminatoire \\
  \bottomrule
\end{tabular}
\end{center}

Cette classification permet à l'utilisateur de filtrer les sports dès le menu principal via une
logique ET combinant les trois critères. Les sports de type \textit{Mixte} apparaissent dans les
résultats des filtres \textit{Points} et \textit{Éliminatoire}.

\subsection{Fonctionnalités par sport}

Les fonctionnalités sont adaptées à la nature de chaque sport. Pour les sports
\textbf{collectifs} (basketball, LoL, football, CS2, volleyball), l'utilisateur peut choisir entre
des statistiques sur les \textbf{équipes} (classement, stats détaillées, bracket, roster) et sur
les \textbf{joueurs} (fiche individuelle, liste paginée). Pour les sports \textbf{individuels}
(tennis, échecs, badminton, StarCraft~II), les fonctionnalités portent directement sur les joueurs.

Le volleyball et le tennis proposent en outre une distinction \textbf{hommes/femmes}, ajoutant
un niveau de navigation supplémentaire. Parmi les fonctionnalités communes à tous les sports :
la \textbf{fiche joueur} (informations personnelles, statistiques clés), les
\textbf{classements} (par victoires, par Elo pour les échecs, par points), les
\textbf{bilans} (victoires/défaites/nuls) et l'\textbf{export CSV} de tout résultat tabulaire.

\subsection{Architecture générale}

L'application s'organise en quatre couches principales. La figure~\ref{fig:archi}
présente cette architecture, dont le diagramme de classes ci-après détaille les relations
orientées objet.

\begin{figure}[h!]
\centering
\begin{tikzpicture}[
  layer/.style={draw, fill=#1, rounded corners=4pt,
                minimum width=10cm, minimum height=1.1cm,
                font=\bfseries\small, align=center},
  arr/.style={->, thick, >=Stealth},
  note/.style={font=\footnotesize\itshape, text width=4.5cm, align=left}
]
  \node[layer=orange!20] (main)    at (0,  5.5) {\_\_main\_\_.py \quad (Interface CLI --- 1\,397 lignes)};
  \node[layer=blue!15]   (analysis)at (0,  3.8) {src/Analysis/ \quad (Modules analytiques par sport)};
  \node[layer=green!15]  (parsers) at (0,  2.1) {src/Parsers/ \quad (Chargeurs CSV)};
  \node[layer=red!15]    (model)   at (0,  0.4) {src/Model/ \quad (Classes métier génériques)};
  \node[layer=gray!15]   (common)  at (0, -1.2) {src/Common/ \quad (Utilitaires partagés)};

  \draw[arr] (main)     -- (analysis);
  \draw[arr] (analysis) -- (parsers);
  \draw[arr] (parsers)  -- (model);
  \draw[arr] (analysis.west) to[out=180, in=180] (common.west);
  \draw[arr] (parsers.east)  to[out=0,   in=0]   (model.east);

  \node[note] at (7.3, 4.65)  {stats, classements, rosters, brackets};
  \node[note] at (7.3, 2.95)  {lecture CSV, nettoyage, jointures};
  \node[note] at (7.3, 1.25)  {Joueur, Match, Equipe, Sport};
  \node[note] at (7.3, -0.35) {décorateurs, helpers};
\end{tikzpicture}
\caption{Architecture en couches de l'application.}
\label{fig:archi}
\end{figure}

\begin{figure}[h!]
\centering
\begin{tikzpicture}[
  cls/.style={draw, fill=white, font=\scriptsize\ttfamily,
              minimum width=2.9cm, align=center},
  enum/.style={draw, fill=Plum!15, font=\scriptsize\ttfamily,
               minimum width=2.4cm, align=center},
  abst/.style={draw, fill=orange!15, font=\scriptsize\ttfamily,
               minimum width=2.9cm, align=center},
  cli/.style={draw, fill=green!10, font=\scriptsize\ttfamily,
              minimum width=2.9cm, align=center},
  inh/.style={-{Triangle[open]}, thick},
  agg/.style={-{Diamond[open]}, thick},
  dep/.style={dashed, -Stealth, thick},
  >=Stealth
]
  % Enums
  \node[enum, minimum height=1.5cm] (tsport) at (-5.5, 6.5) {
    \textbf{TypeSport}\\[2pt]\rule{2.2cm}{0.3pt}\\INDIVIDUEL\\COLLECTIF
  };
  \node[enum, minimum height=1.5cm] (catsport) at (-2.5, 6.5) {
    \textbf{CategorieSport}\\[2pt]\rule{2.2cm}{0.3pt}\\SPORT\\ESPORT
  };
  \node[enum, minimum height=1.8cm] (tcomp) at (0.5, 6.5) {
    \textbf{TypeCompetition}\\[2pt]\rule{2.2cm}{0.3pt}\\POINTS\\ELIMINATOIRE\\MIXTE
  };

  % Sport
  \node[cls, minimum height=2.2cm] (sport) at (-2.5, 3.8) {
    \textbf{Sport} (dataclass frozen)\\[2pt]\rule{2.7cm}{0.3pt}\\
    nom, type\_sport\\categorie, type\_competition\\couleur\\[2pt]
    \rule{2.7cm}{0.3pt}\\
    filtrer(), get\_sport()
  };
  \draw[inh] (sport.north) -- ++(0, 0.4) -| (tsport.south);
  \draw[inh] (sport.north) -- ++(0, 0.4) -| (catsport.south);
  \draw[inh] (sport.north) -- ++(0, 0.4) -| (tcomp.south);

  % StatsSport
  \node[abst, minimum height=1.8cm] (statssport) at (3.5, 3.8) {
    \textbf{StatsSport} (ABC)\\[2pt]\rule{2.7cm}{0.3pt}\\[2pt]
    \rule{2.7cm}{0.3pt}\\
    classement(), roster()\\nb\_joueurs(), nb\_matchs()\\est\_collectif()
  };

  % Joueur / Joueurs
  \node[cls, minimum height=2.4cm] (joueurs) at (-5.5, 0.5) {
    \textbf{Joueur / Joueurs}\\[2pt]\rule{2.7cm}{0.3pt}\\
    col\_id, col\_nom\\col\_naissance, col\_taille\\[2pt]
    \rule{2.7cm}{0.3pt}\\
    get\_par\_nom()\\age\_moyen(), resume()
  };

  % Equipe / Equipes
  \node[cls, minimum height=2.4cm] (equipes) at (-2.5, 0.5) {
    \textbf{Equipe / Equipes}\\[2pt]\rule{2.7cm}{0.3pt}\\
    col\_id, col\_nom\\col\_abrev, col\_ville\\[2pt]
    \rule{2.7cm}{0.3pt}\\
    get\_par\_nom()\\get\_joueurs()
  };

  % Match / Matchs
  \node[cls, minimum height=2.4cm] (matchs) at (0.5, 0.5) {
    \textbf{Match / Matchs}\\[2pt]\rule{2.7cm}{0.3pt}\\
    col\_equipe1/2\\col\_score1/2, col\_gagnant\\[2pt]
    \rule{2.7cm}{0.3pt}\\
    bilan(), classement()\\victoires(), defaites()
  };

  % BaseLoader
  \node[abst, minimum height=2.2cm] (loader) at (-4.0, -3.0) {
    \textbf{BaseLoader} (ABC)\\[2pt]\rule{2.7cm}{0.3pt}\\
    SPORT\_FOLDER\\[2pt]
    \rule{2.7cm}{0.3pt}\\
    \_load\_csv()\\load\_all()
  };

  % Loaders concrets
  \node[cls, minimum height=0.8cm] (bball)  at (-5.8, -5.0) {\textbf{Basketball}};
  \node[cls, minimum height=0.8cm] (tennis) at (-4.0, -5.0) {\textbf{Tennis}};
  \node[cls, minimum height=0.8cm] (lol)    at (-2.2, -5.0) {\textbf{LoL, CS2...}};

  \draw[inh] (bball.north)  -- (loader.south);
  \draw[inh] (tennis.north) -- (loader.south);
  \draw[inh] (lol.north)    -- (loader.south);

  % CLI
  \node[cli, minimum height=2.4cm] (cli) at (3.5, 0.5) {
    \textbf{CLI}\\[2pt]\rule{2.7cm}{0.3pt}\\
    \_back, \_fwd, \_filters\\[2pt]
    \rule{2.7cm}{0.3pt}\\
    run()\\\_page\_accueil()\\\_make\_page\_sport()
  };

  \draw[dep] (cli.west) -- (matchs.east);
  \draw[dep] (loader.east) to[out=0, in=270] (matchs.south);
  \draw[dep] (loader.east) to[out=0, in=270] (joueurs.south);

\end{tikzpicture}
\caption{Diagramme de classes simplifié de l'application.}
\label{fig:classes}
\end{figure}

% ------------------------------------------------------------
\section{Répartition des rôles}
\label{sec:roles}
% ------------------------------------------------------------

Ce projet a suivi plusieurs phases principales. La conception a de loin été la phase la
plus longue, en raison d'un emploi du temps chargé en cette période. Nous avons réfléchi
ensemble à la structure globale de notre future application et il a d'abord fallu comprendre
clairement les attendus du projet.

\textbf{Phase de conception} --- Réflexion collective sur l'architecture, le choix des
sports, la structure des données et l'organisation des modules. Chaque membre a
prototypé ses idées dans un dossier individuel (\fichier{Projet\_Adam},
\fichier{Projet\_Jade}, \fichier{Projet\_Mathias}, \fichier{Projet\_fouad}) avant la fusion
dans le code commun.

\textbf{Phase de développement} --- Les tâches ont été réparties par couche et par sport.
La couche Model (classes génériques Joueur, Equipe, Match) et le BaseLoader ont été
développés en priorité pour fournir les fondations communes. Puis chaque membre a
implémenté les modules d'analyse et les loaders de ses sports attribués. Le module
générique a été développé itérativement au fur et à mesure que les patterns communs
émergeaient.

\textbf{Phase d'intégration} --- Le \fichier{\_\_main\_\_.py} et le moteur CLI ont été
développés pour relier tous les modules. Cette phase a nécessité un travail de
coordination important pour harmoniser les interfaces entre les modules d'analyse.

\textbf{Phase de tests et documentation} --- Les tests unitaires ont été écrits en
parallèle du développement, chaque membre testant ses modules. Le formatage du code
est assuré par black. La documentation est intégrée via les docstrings Python (format NumPy).

\begin{center}
\begin{tabular}{ll}
  \toprule
  \textbf{Membre} & \textbf{Contributions principales} \\
  \midrule
  M.-A. \textsc{Lakhrissi}  & Architecture, Model, CLI, intégration \\
  Jade \textsc{Colin}       & Modélisation, classes de données, tests \\
  Mathias \textsc{Douillard}& Loaders, analyse, données, export \\
  M.~F. \textsc{Chaabane}   & Interface, visualisation, analyse \\
  \bottomrule
\end{tabular}
\end{center}

% ------------------------------------------------------------
\section{Gestion des données}
\label{sec:donnees}
% ------------------------------------------------------------

\subsection{Sources et formats des données}

Les données du projet proviennent de sources publiques variées et sont toutes stockées
au format CSV dans le dossier \fichier{data/}, organisé par sport. Au total, le projet
exploite \textbf{36 fichiers CSV} répartis sur 10 dossiers.

\begin{center}
\begin{tabular}{lcccc}
  \toprule
  \textbf{Sport} & \textbf{Fichiers CSV} & \textbf{Joueurs} & \textbf{Matchs} & \textbf{Source} \\
  \midrule
  Basketball (NBA)      & 3 & 431     & $\sim$1\,200 & NBA API \\
  Football (ligues EU)  & 5 & $\sim$11\,000 & 25\,979 & Open Football \\
  Football Champions L. & 3 & 751     & 125    & Kaggle \\
  Tennis (ATP + WTA)    & 4 & 778     & 5\,765 & Tennis Abstract \\
  League of Legends     & 4 & 50      & 45     & Liquipedia \\
  Échecs (FIDE)         & 2 & 206     & 256    & FIDE \\
  Volleyball (JO 2024)  & 7 & 311     & $\sim$90  & Paris 2024 \\
  CS2                   & 4 & 160     & 106    & Liquipedia \\
  StarCraft~II          & 2 & 32      & 67     & Liquipedia \\
  Badminton (BWF)       & 2 & 88      & 226    & BWF \\
  \midrule
  \textbf{Total}        & \textbf{36} & \textbf{$>$3\,000} & \textbf{$>$30\,000} & \\
  \bottomrule
\end{tabular}
\end{center}

\subsection{Chargement et normalisation (Parsers)}

Le chargement des données est centralisé dans le package \fichier{src/Parsers}. Une
classe abstraite \classe{BaseLoader} fournit l'infrastructure commune : résolution du chemin
vers le dossier de données, lecture CSV via pandas, et conversion automatique des colonnes
de dates. Chaque sport hérite de cette classe en spécifiant simplement son
\code{SPORT\_FOLDER} et en implémentant ses méthodes \code{load\_*()}.

Par exemple, le \classe{BasketballLoader} charge trois fichiers (\fichier{player.csv},
\fichier{team.csv}, \fichier{game.csv}) et retourne un tuple de trois DataFrames.
Le \classe{VolleyballLoader}, plus complexe, charge sept fichiers en séparant joueurs
et coachs par genre. Ce pattern permet d'ajouter un nouveau sport en quelques lignes.

\begin{lstlisting}[caption={Structure minimale d'un loader (exemple Badminton)}]
class BadmintonLoader(BaseLoader):
    SPORT_FOLDER = "badminton"

    def load_all(self):
        players = self._load_csv("player.csv")
        matches = self._load_csv("match.csv", date_cols=["date"])
        return players, matches
\end{lstlisting}

\subsection{Nettoyage et validation}

Le nettoyage des données s'effectue à deux niveaux.

\textbf{Dans les Loaders} : les colonnes de dates sont parsées en \code{datetime} pandas
via \code{pd.to\_datetime(errors='coerce')}, ce qui gère automatiquement les formats
variés et les valeurs manquantes. Pour le Basketball NBA, les tailles au format américain
(\code{"6-8"} pour 6 pieds 8 pouces) sont converties en centimètres :

\begin{lstlisting}[caption={Conversion de taille NBA (BasketballLoader.py)}]
def _height_to_cm(h: str) -> float | None:
    try:
        feet, inches = map(int, h.split("-"))
        return round((feet * 12 + inches) * 2.54, 1)
    except (ValueError, AttributeError):
        return None
\end{lstlisting}

\textbf{Dans les modules d'analyse} : les fonctions appliquent un \code{dropna()} ciblé
sur les colonnes pertinentes avant tout calcul, et utilisent des valeurs de
remplacement \code{"N/A"} pour l'affichage des fiches joueurs. Pour les colonnes
numériques (kills, gold, assists, etc.), les valeurs manquantes sont gérées implicitement
par les méthodes agrégées de pandas.

% ------------------------------------------------------------
\section{Structure du code}
\label{sec:structure}
% ------------------------------------------------------------

\subsection{Organisation des packages}

Le projet suit une organisation en packages Python claire et modulaire, totalisant
environ \textbf{11\,400 lignes de code Python} réparties sur \textbf{69 fichiers}.

\begin{Verbatim}[fontsize=\small, frame=single]
Projet-1A-2026/
|-- __main__.py           # Point d'entree CLI (1 397 lignes)
|-- requirements.txt      # pandas, pytest, pytest-cov, black
|-- README.md
|-- data/                 # 36 fichiers CSV (10 sports)
|-- src/
|   |-- Model/            # Classes generiques
|   |-- Parsers/          # Loaders par sport
|   |-- Analysis/         # Modules analytiques
|   `-- Common/           # Utilitaires partages
|-- test/                 # ~3 800 lignes de tests
|-- Projet_Adam/          # Contributions individuelles
|-- Projet_Jade/
|-- Projet_Mathias/
`-- Projet_fouad/
\end{Verbatim}

Les sous-packages importent rarement les uns des autres directement. Les modules
d'analyse importent uniquement leurs Loaders et le module générique. Le
\fichier{\_\_main\_\_.py} est le seul fichier qui relie tous les modules ensemble,
conformément aux bonnes pratiques de couplage faible.

\subsection{Le fichier \_\_main\_\_.py}

Ce fichier de \textbf{1\,397 lignes} est le chef d'orchestre de l'application.
Il contient deux éléments principaux :

\textbf{La configuration des sports} (\code{\_make\_sports\_config()}) --- un
dictionnaire associant chaque sport à ses fonctionnalités disponibles, organisées par
catégorie (équipe/joueurs) et par genre le cas échéant. Chaque fonctionnalité est décrite
par un label, une fonction d'analyse à appeler, et optionnellement un sélecteur avec
pagination et filtres.

\textbf{La classe CLI} --- implémente un moteur de navigation basé sur des
pages-fonctions (\textit{callables}). Chaque page retourne soit une nouvelle page
(navigation avant), soit \code{'back'}/\code{'forward'}/\code{'quit'} pour la navigation
historique, soit \code{None} pour rester sur place. Ce pattern simple mais puissant
permet de construire des arbres de navigation arbitraires sans machine à états complexe.

\subsection{Le moteur de navigation CLI}

La classe \classe{CLI} gère un \textbf{historique de navigation avec deux piles}
(\code{\_back} et \code{\_fwd}), permettant à l'utilisateur de revenir en arrière ou
d'avancer dans son parcours. Le flux passe par quatre niveaux :

\begin{enumerate}
  \item \textbf{Page d'accueil} --- affiche les sports filtrés et l'option de filtrage.
  \item \textbf{Page sport} --- routeur qui oriente vers la sélection genre,
        équipe/joueurs, ou directement les stats selon la configuration.
  \item \textbf{Page stats} --- liste les fonctionnalités disponibles.
  \item \textbf{Page résultat} --- exécute la fonction d'analyse et propose l'export CSV.
\end{enumerate}

Pour les sélecteurs avec de nombreuses options (joueurs, équipes), un système de
\textbf{pagination par 20 éléments} avec filtres (recherche par nom, catégorie, année
de naissance) permet une navigation fluide même sur de grandes bases (ex. : 443 joueurs
ATP, 11\,000 joueurs dans les ligues européennes de football).

\begin{figure}[h!]
\centering
\begin{tikzpicture}[
  box/.style={draw, rounded corners, fill=NavyBlue!8, font=\small,
              minimum width=3.8cm, minimum height=0.9cm, align=center},
  side/.style={draw, rounded corners, fill=orange!15, font=\small,
               minimum width=3cm, minimum height=0.9cm, align=center},
  arr/.style={->, thick, >=Stealth}
]
  \node[box] (start) at (0, 0)    {Lancement\\(\_\_main\_\_.py)};
  \node[box] (menu1) at (0, -1.8) {Page d'accueil\\(choix du sport + filtres)};
  \node[box] (menu2) at (0, -3.6) {Page sport\\(équipe / joueurs / genre)};
  \node[box] (input) at (0, -5.4) {Sélection paginée\\(20 items, filtres, recherche)};
  \node[box] (result)at (0, -7.2) {Affichage du résultat};
  \node[side] (export)at (5, -7.2) {Export CSV\\(horodaté, utf-8-sig)};
  \node[box] (retour)at (0, -9.0) {Retour au menu};

  \draw[arr] (start)  -- (menu1);
  \draw[arr] (menu1)  -- (menu2);
  \draw[arr] (menu2)  -- (input);
  \draw[arr] (input)  -- (result);
  \draw[arr] (result) -- node[above, font=\footnotesize] {[e]} (export);
  \draw[arr] (result) -- (retour);
  \draw[arr] (retour.west) to[out=180, in=180, distance=2.5cm]
    node[left, font=\footnotesize] {[b] retour} (menu1.west);
\end{tikzpicture}
\caption{Flux de navigation de l'interface CLI.}
\label{fig:flux}
\end{figure}

\subsection{Le module générique (générique.py)}

Ce module de \textbf{733 lignes} mutualise les fonctionnalités communes à tous les
sports. Il fournit cinq fonctions principales :

\begin{itemize}
  \item \code{formater\_roster} --- formate le roster d'une équipe avec les colonnes
        Rôle, Nom complet, Pseudo (e-sport uniquement), Nationalité, Date de naissance.
        Gère les coachs en tête de liste et les données manquantes.
  \item \code{afficher\_bracket} --- construit et affiche un tableau de tournoi
        éliminatoire en ASCII, avec les gagnants mis en évidence en vert (ANSI).
        Supporte les compétitions à aller-retour.
  \item \code{afficher\_classement} --- produit un classement par points (3 pts victoire,
        1 pt nul, 0 pt défaite) avec séparateur de qualification optionnel.
  \item \code{fiche\_joueur} --- génère une fiche complète d'un joueur adaptée à chaque
        sport via un dictionnaire de labels.
  \item \code{lister\_joueurs} --- liste les joueurs d'une équipe ou d'un pays avec
        filtrage et tri alphabétique.
\end{itemize}

% ------------------------------------------------------------
\section{Programmation orientée objet}
\label{sec:poo}
% ------------------------------------------------------------

Le projet fait un usage mesuré mais pertinent de la programmation orientée objet. Comme
le souligne le README, les classes sont principalement utiles pour « passer des données
formatées à travers l'application » ; l'essentiel du traitement de données s'appuie sur
les DataFrames pandas. Les classes generiques avec leur \textbf{mapping de colonnes
paramétrable} offrent un bon compromis entre abstraction et simplicité.

\subsection{La classe Sport et les enums de classification}

La classe \classe{Sport}, définie comme un \code{@dataclass(frozen=True)}, encapsule
les métadonnées d'un sport : nom, type (collectif/individuel), catégorie (sport/e-sport),
type de compétition (points/éliminatoire/mixte), et couleur d'interface. Trois enums
structurent ces axes :

\begin{lstlisting}[caption={Enums de classification des sports (src/Model/sport.py)}]
class TypeSport(Enum):
    INDIVIDUEL = "individuel"
    COLLECTIF  = "collectif"

class CategorieSport(Enum):
    SPORT  = "sport"
    ESPORT = "e-sport"

class TypeCompetition(Enum):
    POINTS       = "points"
    ELIMINATOIRE = "eliminatoire"
    MIXTE        = "mixte"
\end{lstlisting}

Le registre \code{SPORTS} (liste de 9 instances) et la fonction \code{filtrer()} permettent
de combiner les critères en logique ET pour le menu de filtrage du CLI. Des propriétés
dérivées comme \code{est\_en\_points} gèrent la logique des sports de type Mixte.

\subsection{Les classes Joueur et Joueurs}

La classe \classe{Joueur} enveloppe une \code{pd.Series} (ligne de DataFrame) avec un
mapping de colonnes paramétrable : \code{col\_id}, \code{col\_nom},
\code{col\_naissance}, \code{col\_equipe}, \code{col\_taille}. Des propriétés calculées
comme \code{age} dérivent automatiquement de la date de naissance via
\code{pd.Timestamp.today()}.

La classe \classe{Joueurs} (collection) wrappe un DataFrame entier et fournit des méthodes
de filtrage (\code{get\_par\_nom}, \code{get\_par\_equipe}, \code{filtrer}) et de
statistiques (\code{age\_moyen}, \code{taille\_moyenne}, \code{resume}). Elle implémente
les protocoles Python standard (\code{\_\_len\_\_}, \code{\_\_iter\_\_},
\code{\_\_getitem\_\_}) pour une utilisation naturelle.

La compatibilité multi-sport est garantie par le paramétrage des colonnes à
l'initialisation :

\begin{center}
\small
\begin{tabular}{lllll}
  \toprule
  \textbf{Sport} & \textbf{col\_id} & \textbf{col\_nom} & \textbf{col\_equipe} & \textbf{col\_taille} \\
  \midrule
  Basketball & \code{person\_id} & \code{full\_name} & \code{team\_id} & \code{height\_cm} \\
  Tennis     & \code{player\_id} & \code{full\_name} & ---             & \code{height} \\
  LoL        & \code{pseudo}     & \code{name}       & \code{team}     & --- \\
  Volleyball & \code{name}       & \code{name}       & \code{country\_code} & \code{height} \\
  \bottomrule
\end{tabular}
\end{center}

\subsection{Les classes Equipe et Equipes}

Suivant le même pattern que \classe{Joueur}/\classe{Joueurs}, ces classes modélisent
les équipes avec les colonnes \code{col\_id}, \code{col\_nom}, \code{col\_abrev},
\code{col\_ville}. La méthode \code{get\_joueurs(joueurs: Joueurs)} permet de récupérer
les joueurs d'une équipe en croisant les deux collections --- illustrant la composition
orientée objet.

\subsection{Les classes Match et Matchs}

La classe \classe{Match} représente un match avec les colonnes d'équipes, de scores et
de gagnant. La propriété \code{gagnant} implémente une logique à deux niveaux : elle
utilise d'abord la colonne dédiée si elle existe, sinon déduit le gagnant depuis les scores.
La collection \classe{Matchs} fournit des statistiques agrégées : \code{bilan},
\code{classement}, \code{score\_moyen}, avec filtrage par équipe, date et saison.

\subsection{La classe abstraite StatsSport}

Cette ABC définit le contrat commun pour tous les modules analytiques :
\code{classement(n)}, \code{nb\_joueurs()}, \code{nb\_matchs()},
\code{stats\_entite(nom)}. La méthode optionnelle \code{roster()} est disponible pour
les sports collectifs. La méthode \code{est\_collectif()} détecte automatiquement si
\code{roster} a été surchargée via l'introspection Python.

\subsection{Le pattern BaseLoader}

\classe{BaseLoader} est une ABC qui centralise la mécanique de chargement CSV.
Elle implémente le \textbf{patron de conception Template Method} : les sous-classes n'ont
qu'à déclarer \code{SPORT\_FOLDER} et à implémenter les méthodes \code{load\_*()},
tout le reste étant géré par la classe mère.

\begin{lstlisting}[caption={BaseLoader.py -- infrastructure commune à tous les loaders}]
class BaseLoader(ABC):
    _DATA_ROOT: Path = Path(__file__).parent.parent.parent / "data"
    SPORT_FOLDER: str = ""

    @property
    def root(self) -> Path:
        return self._DATA_ROOT / self.SPORT_FOLDER

    def _load_csv(self, filename, dtype=None, date_cols=None):
        df = pd.read_csv(self.root / filename, dtype=dtype or {})
        for col in (date_cols or []):
            if col in df.columns:
                df[col] = pd.to_datetime(df[col], errors="coerce")
        return df
\end{lstlisting}

Les modules d'analyse utilisent également un \textbf{pattern singleton de chargement
paresseux} : les données ne sont chargées qu'au premier appel et mises en cache pour
les appels suivants, garantissant de bonnes performances.

% ------------------------------------------------------------
\section{Tests et validation}
\label{sec:tests}
% ------------------------------------------------------------

\subsection{Stratégie de test}

Le projet adopte une stratégie de tests unitaires rigoureuse avec pytest. Conformément
aux recommandations du README, les tests n'utilisent jamais les données réelles (fichiers
CSV) mais des \textbf{DataFrames synthétiques} créés dans des fixtures pytest. Cela
garantit des tests rapides, reproductibles et indépendants de l'environnement de données.

Les modules de test sont organisés en miroir du code source. Chaque module d'analyse
sport dispose d'un \code{monkeypatch} des variables globales (\code{\_players},
\code{\_teams}, \code{\_matches}) avec des données synthétiques via les fixtures.
Chaque test vérifie non seulement le résultat attendu, mais aussi les cas d'erreur
(\code{ValueError} pour les entrées invalides, gestion des données manquantes, etc.).

\subsection{Couverture}

\begin{center}
\begin{tabular}{lcc}
  \toprule
  \textbf{Package} & \textbf{Fichiers de test} & \textbf{Lignes de test} \\
  \midrule
  Model    & 5  & 984   \\
  Analysis & 10 & 2\,188 \\
  Parsers  & 1  & 531   \\
  Common   & 1  & 93    \\
  \midrule
  \textbf{Total} & \textbf{17} & \textbf{3\,796} \\
  \bottomrule
\end{tabular}
\end{center}

Les tests se lancent via :

\begin{lstlisting}[language=bash, numbers=none]
python -m pytest --cov=src --cov-report=term-missing
\end{lstlisting}

\subsection{Formatage et linting}

L'ensemble du code est formaté avec \textbf{black} (version 25.1.0) et vérifié avec
\textbf{Flake8}. Le respect de ces outils garantit une cohérence stylistique sur
l'ensemble du projet et est contrôlé systématiquement avant chaque merge.

% ------------------------------------------------------------
\section{Résultats}
\label{sec:resultats}
% ------------------------------------------------------------

L'application finale permet d'explorer les statistiques de neuf sports à travers une interface
CLI complète et fonctionnelle. Nous présentons ci-dessous cinq scénarios d'utilisation
représentatifs, illustrés par des captures d'écran.

\subsection{Scénario 1 : Navigation et filtrage des sports}

\begin{scenario}[Scénario 1 --- Navigation et filtrage]
L'utilisateur lance l'application et arrive sur le menu principal listant les neuf sports.
Il active le filtre \textit{Collectif} + \textit{E-sport} : seuls League of Legends et
CS2 apparaissent. Il sélectionne CS2, puis \textit{Équipe}, puis \textit{Classement} pour
obtenir le classement de la phase de groupes, affiché sous forme de tableau trié
(équipe, victoires, défaites, matchs joués, points).
\end{scenario}

\screenshot{Menu principal avec filtres actifs (Collectif + E-sport)}{
  Seuls League of Legends et CS2 sont affichés après application du filtre.
  La barre de filtres active est visible en en-tête.
}{menu_filtres.png}

\subsection{Scénario 2 : Fiche joueur avec sélecteur paginé}

\begin{scenario}[Scénario 2 --- Fiche joueur avec sélecteur paginé]
L'utilisateur choisit Tennis > Hommes > Fiche d'un joueur. Le sélecteur paginé affiche les
443 joueurs ATP par pages de 20. Il utilise le filtre \textit{Pays (IOC)} pour ne voir que
les joueurs français, puis sélectionne un joueur. La fiche affiche le nom complet, la
nationalité, la date de naissance, la taille, le nombre de victoires et défaites, et les
performances par surface.
\end{scenario}

\screenshot{Sélecteur paginé avec filtre « Pays : FRA » actif}{
  Liste des joueurs ATP français (20 par page).
  Navigation : [<] Préc. | [>] Suiv. | [f] Filtrer | [r] Reset.
}{filtre_tennis_francais.png}

\screenshot{Fiche individuelle d'un joueur ATP}{
  Informations complètes : nom, nationalité, date de naissance, taille,
  bilan victoires/défaites, statistiques par surface.
  Option [e] Exporter en CSV disponible.
}{fiche_moutet.png}

\subsection{Scénario 3 : Bracket de tournoi}

\begin{scenario}[Scénario 3 --- Bracket de tournoi]
L'utilisateur sélectionne Football Champions League > Équipe > Bracket. L'application
affiche l'arbre du tournoi avec les résultats de chaque tour éliminatoire, reconstruit à
partir des données de matchs. Les gagnants sont mis en évidence en vert. Ce bracket est
également disponible pour le tennis (Grand Chelem), le badminton et CS2.
\end{scenario}

\screenshot{Bracket éliminatoire de la Ligue des Champions UEFA}{
  Arbre du tournoi affiché en ASCII. Les équipes gagnantes apparaissent
  en vert (codes ANSI). Chaque colonne représente un tour (R16, QF, SF, F).
}{bracket_champions.png}

\subsection{Scénario 4 : Statistiques d'équipe e-sport}

\begin{scenario}[Scénario 4 --- Statistiques d'équipe LoL]
L'utilisateur choisit League of Legends > Équipe > Stats d'une équipe. Après sélection
de G2 Esports dans le sélecteur paginé, l'application affiche les moyennes par match :
kills, assists, morts, or total, tourelles, dragons, barons et durée moyenne des parties
en minutes.
\end{scenario}

\screenshot{Statistiques moyennes par match de G2 Esports (EMEA LEC)}{
  Tableau de stats : kills/assists/morts/or/tourelles/dragons/barons.
  Option [e] disponible pour export CSV.
}{stats_g2.png}

\subsection{Scénario 5 : Export CSV}

\begin{scenario}[Scénario 5 --- Export CSV]
Après l'affichage de n'importe quel résultat tabulaire (classement, roster, stats, liste
de joueurs), l'utilisateur appuie sur \code{[e]} pour exporter le DataFrame en CSV.
Un nom de fichier horodaté est proposé par défaut
(\code{\{label\}\_YYYYMMDD\_HHMMSS.csv}), encodé en UTF-8-BOM pour compatibilité
Excel, et sauvegardé dans le répertoire courant.
\end{scenario}

\screenshot{Confirmation d'export CSV après un classement NBA}{
  Message de confirmation avec le chemin du fichier généré.
  Exemple : \texttt{classement\_nba\_20260508\_143022.csv}
}{export_csv.png}

% ------------------------------------------------------------
\section{Conclusion}
\label{sec:conclusion}
% ------------------------------------------------------------

Ce projet a été l'occasion de mettre en pratique les concepts fondamentaux de la
programmation Python dans un contexte concret et stimulant. La construction d'une
application multi-sports nous a confrontés à des défis techniques variés : concevoir des
classes suffisamment génériques pour s'adapter à neuf sports aux structures de données
très différentes, normaliser des sources hétérogènes, et maintenir un code modulaire et
testable à mesure que le projet grossissait.

Le choix d'une architecture en couches (Parsers $\rightarrow$ Model $\rightarrow$
Analysis $\rightarrow$ CLI) s'est révélé particulièrement pertinent. L'ajout d'un nouveau
sport se résume à créer un Loader, un module d'analyse, et une entrée dans la
configuration du \fichier{\_\_main\_\_.py} --- sans toucher au moteur de navigation ni
aux classes génériques. Cette extensibilité nous a permis de passer de cinq à neuf sports
sans refactoring majeur.

L'utilisation parcimonieuse des classes, combinée à la puissance de pandas pour le
traitement tabulaire, correspond à une approche pragmatique adaptée au contexte du
projet. Les classes génériques (\classe{Joueur}, \classe{Equipe}, \classe{Match}) avec
leur mapping de colonnes paramétrable offrent un bon compromis entre abstraction et
simplicité, comme en témoignent les 3\,796 lignes de tests garantissant la fiabilité
des composants critiques.

Parmi les améliorations possibles, on pourrait envisager : l'ajout d'une interface
graphique (web ou terminal enrichi avec \textit{rich}/\textit{textual}), l'intégration
de visualisations statistiques (matplotlib/seaborn), des fonctionnalités d'analyse
comparative entre sports, ou encore la connexion à des API en temps réel. Ce projet nous
a en tout cas permis de développer une solide expérience en architecture logicielle,
en travail collaboratif via Git, et en exploitation de données structurées avec Python.

% ------------------------------------------------------------
\newpage
\begin{thebibliography}{9}
  \bibitem{pandas}
    McKinney, W. (2010).
    \textit{Data Structures for Statistical Computing in Python}.
    Proceedings of the 9th Python in Science Conference.

  \bibitem{pep8}
    van Rossum, G., Warsaw, B., Coghlan, A. (2001).
    \textit{PEP 8 -- Style Guide for Python Code}.
    Python Software Foundation.

  \bibitem{pytest}
    Krekel, H. et al. (2004--2024).
    \textit{pytest: helps you write better programs}.
    \url{https://docs.pytest.org/}

  \bibitem{kaggle_nba}
    Jeu de données NBA 2022--2023.
    Kaggle (source publique).

  \bibitem{tennis_abstract}
    Sackmann, J.
    \textit{Tennis Abstract --- ATP/WTA match results}.
    \url{https://www.tennisabstract.com/}

  \bibitem{liquipedia}
    Liquipedia --- base de données e-sport (LoL, CS2, StarCraft~II).
    \url{https://liquipedia.net/}
\end{thebibliography}

% ============================================================
% ANNEXES
% ============================================================
\newpage
\appendix
\renewcommand{\thesection}{\Alph{section}}
\section*{Annexes}
\addcontentsline{toc}{section}{Annexes}

% -----------------------------------------------
\subsection*{Annexe A --- Bilan personnel de Jade Colin}
\addcontentsline{toc}{subsection}{Annexe A --- Bilan personnel de Jade Colin}

\begin{bilanbox}
J'ai personnellement trouvé ce projet assez difficile. Au début, j'ai eu besoin de temps
pour bien comprendre les attentes ainsi que la manière de démarrer le projet. Étant débutante
en programmation, certaines consignes me paraissaient complexes et il n'était pas évident de
savoir par où commencer.

\medskip
De plus, notre groupe était composé de personnes qui ne se connaissaient pas auparavant, ce
qui a rendu la cohésion et l'organisation plus difficiles au départ. Malgré cela, nous avons
progressivement réussi à communiquer, à échanger nos idées et à trouver une méthode de travail
efficace.

\medskip
Après plusieurs discussions et une meilleure compréhension des objectifs du projet, j'ai
commencé à mieux appréhender les différentes étapes de réalisation. Ce projet m'a notamment
permis d'approfondir mes connaissances en programmation orientée objet, en apprenant à
concevoir des classes et des objets répondant à un besoin concret.

\medskip
L'une des principales difficultés a été de développer un code générique, capable de fonctionner
avec plusieurs bases de données différentes, ce qui demandait une réflexion plus poussée sur
l'architecture du programme et sa modularité.

\medskip
Au final, même si ce projet a été exigeant et parfois difficile, il m'a permis de progresser,
de mieux comprendre les principes de la programmation orientée objet et d'améliorer mes
compétences techniques ainsi que ma capacité à travailler en équipe.
\end{bilanbox}

% -----------------------------------------------
\subsection*{Annexe B --- Bilan personnel de Mathias Douillard}
\addcontentsline{toc}{subsection}{Annexe B --- Bilan personnel de Mathias Douillard}

\begin{bilanbox}
Ce projet a été source de difficultés tout au long de sa confection. J'ai trouvé que nous
avions très peu d'informations à notre disposition au début du projet, ce qui provient
probablement de l'absence de base de données initiale. Il était compliqué de se projeter sur
le fonctionnement concret d'une application car on ne savait pas exactement ce qu'elle ferait.

\medskip
Avec du recul, je comprends que l'objectif derrière cela était que l'on se penche uniquement
sur l'architecture du projet et la manière de réfléchir à la création d'une application.
Cependant, cela n'était pas assez clair et il était assez compliqué de bien comprendre les
attendus de ce projet. De plus, les informations au sujet des bases de données ainsi que du
projet arrivaient au compte-gouttes --- que ce soit les attendus ou les barèmes de notation.
Cela nous a forcés à modifier des éléments que nous avions déjà conçus.

\medskip
Enfin, je trouve que les attendus lors des premières séances en autonomie étaient assez élevés.
Je sors personnellement de prépa D2 et je découvre la programmation depuis ce début d'année.
En dehors de l'UE de remise à niveau et de programmation orientée objet, je n'avais jamais
codé de ma vie. J'ai donc eu énormément de mal à comprendre ne serait-ce que les notions
évoquées lors de ces séances. Combiné aux éléments indiqués plus haut, j'ai rencontré
énormément de difficultés tout au long de ce projet.

\medskip
Malgré tout, je retire de bonnes choses de ce projet. Je considère avoir pris conscience de
l'étendue presque infinie des possibilités liées au code. Je ne m'attendais pas à mettre au
point une application au cours de ma première année à l'ENSAI et j'en suis satisfait, même si
elle n'est évidemment pas parfaite. Je comprends désormais beaucoup mieux le fonctionnement
et la structure du code.
\end{bilanbox}

% -----------------------------------------------
\subsection*{Annexe C --- Bilan personnel de Mohammed-Adam Lakhrissi}
\addcontentsline{toc}{subsection}{Annexe C --- Bilan personnel de Mohammed-Adam Lakhrissi}

\begin{bilanbox}
Ce projet m'a permis de découvrir les réalités du développement logiciel~: gestion des
données, structuration du code par modules, importance des tests. L'utilisation de datasets
réels m'a fait prendre conscience que la qualité des données et la cohérence de leur
organisation est souvent le premier obstacle à tout travail d'analyse, opéré avant même
d'écrire la première ligne de code.

\medskip
J'ai pris également conscience que l'on ne peut pas se passer d'une bonne organisation de son
code dès le départ~: une mauvaise architecture en tout début est souvent un problème plus
difficile à corriger par la suite. De même, le développement de tests peut se révéler trop
court, ce qui donne l'illusion que tout fonctionne alors que des bugs plus subtils peuvent
persister.

\medskip
Plus personnellement, ce projet m'a appris également à mieux décomposer un problème compliqué
en sous-tâches, et à ne pas sous-estimer le temps à allouer à la phase de débogage. C'est
souvent là où l'on apprend le plus.
\end{bilanbox}

% -----------------------------------------------
\subsection*{Annexe D --- Bilan personnel de Mohamed Fouad Chaabane}
\addcontentsline{toc}{subsection}{Annexe D --- Bilan personnel de Mohamed Fouad Chaabane}

\begin{bilanbox}
Ce projet a été pour moi une expérience à la fois stimulante et parfois frustrante, ce qui
est sans doute le signe qu'on a vraiment appris quelque chose.

\medskip
Au départ, j'avoue avoir sous-estimé la difficulté. On s'est dit~: «~on charge des CSV, on
fait quelques stats, c'est bon.~» La réalité a été bien différente. Le premier vrai obstacle
est arrivé dès la phase de conception~: comment faire travailler ensemble quatre personnes sur
quatre sous-projets sans que tout se casse à l'intégration~? Ce qui m'a le plus marqué, c'est
la valeur concrète de la programmation orientée objet. Au début, ça ressemble à beaucoup de
code pour un résultat qui semble identique à ce qu'on aurait fait avec des fonctions simples.

\medskip
J'ai aussi appris l'importance des tests. Si je devais refaire ce projet, je passerais deux
fois plus de temps sur la conception avant d'écrire la première ligne de code. On aurait dû
définir ensemble, dès le début, les interfaces entre modules --- pas les classes entières,
juste les signatures des méthodes importantes et les formats de données échangées. Ça nous
aurait évité les heures perdues à déboguer des imports cassés ou des attributs qui n'existaient
pas sous le bon nom.

\medskip
Ce que je retiens surtout, c'est que le travail en équipe sur du code demande presque autant
de communication que de programmation. Les outils Git, les classes abstraites et les tests ne
font que soutenir cette communication. Ils ne la remplacent pas.
\end{bilanbox}

\end{document}
